\documentclass[conference]{IEEEtran}
\IEEEoverridecommandlockouts

% Pacotes essenciais para o idioma português e matemática
\usepackage[utf8]{inputenc}
\usepackage[T1]{fontenc}
% \usepackage[brazil]{babel}
\usepackage[brazilian]{babel}
\usepackage{cite}
\usepackage{amsmath,amssymb,amsfonts}
\usepackage{algorithmic}
\usepackage{graphicx}
\usepackage{textcomp}
\usepackage{xcolor}
\usepackage{booktabs}

\begin{document}

\title{Predição de Óbito ou Readmissão em Pacientes com Insuficiência Cardíaca: Uma Abordagem de Aprendizado de Máquina}

\author{\IEEEauthorblockN{[Nome do Colega]}
\IEEEauthorblockA{\textit{Dep. de Engenharia Elétrica} \\
\textit{Universidade Federal de Minas Gerais (UFMG)}\\
Belo Horizonte, Brasil \\
email.do.colega@dominio.com}
\and
\IEEEauthorblockN{Luis Henrique Cordeiro Rodrigues}
\IEEEauthorblockA{\textit{Dep. de Engenharia Elétrica} \\
\textit{Universidade Federal de Minas Gerais (UFMG)}\\
Belo Horizonte, Brasil \\
luishenrique.29@hotmail.com}
}

\maketitle

\begin{abstract}
A insuficiência cardíaca (IC) é uma das principais causas de hospitalização e apresenta altas taxas de morbidade e mortalidade pós-alta. Este estudo tem como objetivo desenvolver um modelo de aprendizado de máquina capaz de prever o desfecho combinado de óbito ou readmissão hospitalar em até 6 meses após a internação de pacientes com IC. Utilizando uma base de dados clínicos reais com 2.008 pacientes e atributos demográficos, sinais vitais, comorbidades e exames laboratoriais, o trabalho detalha as etapas de análise exploratória, engenharia de atributos e tratamento de dados ausentes com padrão não-aleatório (NMAR).
\end{abstract}

\begin{IEEEkeywords}
Insuficiência Cardíaca, Aprendizado de Máquina, Prontuários Eletrônicos, Prognóstico Clínico, Bioengenharia, Classificação Binária.
\end{IEEEkeywords}

\section{Introdução}
A insuficiência cardíaca (IC) afeta milhões de pessoas globalmente e representa um desafio significativo para os sistemas de saúde devido à sua elevada taxa de readmissão hospitalar e mortalidade associada a curto e médio prazo \cite{b3}. Identificar preditores precoces de agravamento clínico durante a internação é fundamental para a transição segura dos cuidados e a otimização de recursos hospitalares.

Com o advento dos Registros Eletrônicos de Saúde (EHR, do inglês \textit{Electronic Healthcare Records}), enormes volumes de dados passaram a ser rotineiramente coletados, desde sinais vitais na admissão até painéis laboratoriais complexos e ecocardiogramas \cite{b2}. A aplicação de técnicas de aprendizagem estatística e inteligência computacional sobre esses registros permite extrair padrões multivariados não lineares, auxiliando a prática médica com modelos de suporte à decisão clínica \cite{b4}.

O presente trabalho propõe o desenvolvimento e a avaliação comparativa de algoritmos de classificação supervisionada para prever o risco de um paciente com IC evoluir para óbito ou necessitar de readmissão hospitalar em uma janela de até 6 meses. O estudo documenta o fluxo de processamento de dados, com ênfase na análise exploratória criteriosa, na prevenção de vazamento de informações (\textit{data leakage}) e na modelagem metodológica de dados faltantes estruturais inerentes à rotina clínica.

\section{Metodologia}

\subsection{Base de Dados e Formulação do Alvo}
Os dados utilizados neste estudo são provenientes do repositório \textit{PhysioNet} (versão 1.3) \cite{b1}, consistindo em uma coorte observacional retrospectiva de 2.008 pacientes consecutivamente hospitalizados com diagnóstico de insuficiência cardíaca no \textit{Zigong Fourth People's Hospital} (Sichuan, China) entre dezembro de 2016 e junho de 2019 \cite{b2}. O critério de inclusão baseou-se nas diretrizes da Sociedade Europeia de Cardiologia (ESC), englobando sinais clínicos típicos e elevação de peptídeos natriuréticos (BNP/NT-proBNP) \cite{b3}.

O conjunto de dados original compreende 166 variáveis clínicas distribuídas entre informações demográficas, sinais vitais na admissão, parâmetros ecocardiográficos, escalas de gravidade funcional (classificação NYHA, grau de Killip e Índice de Comorbidade de Charlson), exames laboratoriais realizados no primeiro dia de internação e o registro diário de prescrição farmacológica \cite{b1, b2}.

A variável resposta (\textit{target}) foi estruturada como um problema de classificação binária combinando o \textbf{óbito ou a readmissão hospitalar em até 6 meses} \cite{b4}. A adoção desse desfecho composto mitiga o viés de risco competitivo da readmissão isolada, no qual pacientes que evoluem para óbito precocemente após a alta seriam erroneamente rotulados como casos sem evento. Na coorte analisada, a prevalência do desfecho positivo foi de 41,5\% (833 pacientes, sendo 776 readmissões e 57 óbitos), configurando um equilíbrio de classes adequado para o treinamento supervisionado. Para prevenir o vazamento de dados (\textit{data leakage}), todas as 14 variáveis referentes a informações de alta e acompanhamento pós-internação foram rigorosamente descartadas da matriz de preditores antes de qualquer transformação estatística.

\subsection{Análise Exploratória e Características Basais}
A análise exploratória de dados (EDA) foi conduzida para investigar a distribuição das covariáveis, a completude dos registros e a presença de observações atípicas (\textit{outliers}). O mapeamento da taxa de ausência revelou que a falta de preenchimento nos exames laboratoriais não ocorre de forma puramente aleatória, mas em blocos consistentes com a não solicitação de painéis diagnósticos inteiros (padrão \textit{Not Missing At Random} --- NMAR), destacando-se a gasometria arterial (50,5\% de ausência média) e o perfil lipídico (50,7\%). Na prática clínica, a omissão desses painéis reflete a avaliação médica de estabilidade hemodinâmica ou metabólica do paciente na admissão, carregando sinal prognóstico relevante.

A Tabela~\ref{tab:caracteristicas_basais} apresenta as características basais estratificadas entre o grupo estável e o grupo que atingiu o desfecho negativo. Observa-se maior concentração de pacientes nas classes funcionais mais graves da escala NYHA (classe IV: 36,7\% contra 26,4\%) e no grau IV de Killip (5,2\% contra 1,4\%) entre os indivíduos que evoluíram para óbito ou readmissão, além de menor média de pressão arterial sistólica na admissão. Adicionalmente, a inspeção de marcadores laboratoriais como o peptídeo natriurético cerebral (BNP) evidenciou forte assimetria à direita e valores extremos, fundamentando a adoção de estimadores robustos (baseados na mediana) e normalização de escala nas etapas subsequentes de engenharia de atributos.

\begin{table}[htbp]
\caption{Características Basais Estratificadas pelo Desfecho}
\begin{center}
\begin{tabular}{|l|c|c|}
\hline
\textbf{Característica} & \textbf{Estável (0)} & \textbf{Óbito/Readmissão (1)} \\
\hline
\hline
\multicolumn{3}{|c|}{\textit{Variáveis Contínuas/Discretas (Média $\pm$ DP)}} \\
\hline
IMC ($\text{kg/m}^2$) & $21.88 \pm 13.54$ & $21.66 \pm 13.81$ \\
Pressão Sistólica (mmHg) & $132.93 \pm 24.42$ & $128.42 \pm 24.97$ \\
Fração de Ejeção --- LVEF (\%) & $50.79 \pm 13.49$ & $50.54 \pm 12.86$ \\
Índice de Charlson (CCI) & $1.79 \pm 0.89$ & $1.96 \pm 1.05$ \\
\hline
\multicolumn{3}{|c|}{\textit{Variáveis Categóricas (\%)}} \\
\hline
Sexo Masculino & 40.8\% & 43.9\% \\
Classe NYHA III & 53.1\% & 49.8\% \\
Classe NYHA IV & 26.4\% & 36.7\% \\
Killip Grau III & 19.8\% & 19.1\% \\
Killip Grau IV & 1.4\% & 5.2\% \\
\hline
\end{tabular}
\label{tab:caracteristicas_basais}
\end{center}
\end{table}

\begin{thebibliography}{00}

\bibitem{b1} Z. Zhang, L. Cao, Y. Zhao, Z. Xu, R. Chen, L. Lv, and P. Xu, ``Hospitalized patients with heart failure: integrating electronic healthcare records and external outcome data,'' PhysioNet, version 1.3, May 2022. [Online]. Available: https://doi.org/10.13026/5m60-vs44

\bibitem{b2} Z. Zhang, L. Cao, R. Chen, Y. Zhao, L. Lv, Z. Xu, and P. Xu, ``Electronic healthcare records and external outcome data for hospitalized patients with heart failure,'' \textit{Scientific Data}, vol. 8, no. 1, p. 46, Feb. 2021.

\bibitem{b3} P. Ponikowski \textit{et al.}, ``2016 ESC Guidelines for the diagnosis and treatment of acute and chronic heart failure,'' \textit{European Heart Journal}, vol. 37, no. 27, pp. 2129--2200, Jul. 2016.

\bibitem{b4} G. James, D. Witten, T. Hastie, R. Tibshirani, and J. Taylor, \textit{An Introduction to Statistical Learning with Applications in Python}. Cham, Switzerland: Springer, 2023.

\end{thebibliography}

\end{document}